\subsection*{Riesgos legales}
\addcontentsline{toc}{subsection}{Riesgos legales}

Los riesgos legales hablan sobre las posibles consecuencias legales que la compañía podría afrontar como efecto de la toma o desconocimiento de acciones que se interpretan como ilegales o inapropiadas. Estos riesgos se manifiestan en varias situaciones, como por ejemplo las relaciones laborales, cumplimiento de regulaciones vigentes, problemas con los derechos de propiedad intelectual, problemas fiscales, seguridad de la información y muchas otras áreas más.

Las repercusiones de estos escenarios pueden generar la amenaza de demandas civiles, imposición de multas o sanciones por parte de las entidades gubernamentales, e inclusive, en el peor de los casos, poder enfrentar cargos penales. 

Para nuestra plataforma, se identifican los siguientes riesgos:

\begin{enumerate}
	\item \textbf{Protección de datos personales:} Debido a que es un servicio de planificación financiera, la empresa va a manejar información personal financiera de usuarios que podría ser altamente sensible. Además de métodos de pago con el tema de las suscripciones. Por tanto, para mitigar estos riesgos, se debe asegurar el cumplimiento de la protección de datos, la realización de auditorías regulares, establecer los consentimientos de la plataforma con el usuario, y establecer protocolos de seguridad informática robustos.
	
	\item \textbf{Cumplimiento normativo:} Debido a ser una plataforma edtech con ciertas características de una fintech, estamos sujetos a seguir los registros dados por la Superintendencia Financiera y cumplir con las regulaciones emergentes del ecosistema fintech, para moderar estos riesgos, nos aseguraremos de implementar correctamente los registros ante las respectivas autoridades pertinentes, además de consultar a la Superintendencia y otras entidades en caso de dudas sobre la autorización, vigilancia o registro en específico de cierto tipo de actividades.
	
	\item \textbf{Derechos de propiedad:} Implicación de posibles conflictos generados por el manejo inadecuado de la protección de derechos de autor de los desarrollos tecnológicos propios (como el contenido educativo y los algoritmos de machine learning) además de la implementación de derechos de terceros. Por esto, debemos establecer políticas claras para no alterar aquellos derechos, entregar créditos a los respectivos autores y consultas sobre registro de derechos de autor ante la Superintendencia.
\end{enumerate}
